% -----------------------------------------------------------------------------
% DEXPI Beamer Theme
% -----------------------------------------------------------------------------
% Copyright (c) 2026 DEXPI e.V.
%
% This work is licensed under the Creative Commons Attribution-ShareAlike 4.0
% International License (CC BY-SA 4.0).
%
% You are free to:
%   - Share: copy and redistribute the material in any medium or format
%   - Adapt: remix, transform, and build upon the material for any purpose
%
% Under the following terms:
%   - Attribution: You must give appropriate credit to DEXPI e.V.,
%     provide a link to the license, and indicate if changes were made.
%   - ShareAlike: If you remix, transform, or build upon the material,
%     you must distribute your contributions under the same license.
%
% Full license text:
%   https://creativecommons.org/licenses/by-sa/4.0/
%
%  Author: Michael Wiedau
% -----------------------------------------------------------------------------

%----------------------------------------------------------------------------------------
%    PACKAGES AND THEMES
%----------------------------------------------------------------------------------------

\documentclass[aspectratio=169,xcolor=dvipsnames]{beamer}
\usetheme{DEXPI}

\usepackage{hyperref}
\usepackage{graphicx}
\usepackage{booktabs}
\usepackage{pgfgantt}


%----------------------------------------------------------------------------------------
%    TITLE PAGE
%----------------------------------------------------------------------------------------

\title{\LaTeX-Beamer Theme for DEXPI}
\subtitle{Licensed under CC-BY-SA}

\author{Submitted by Dr. Michael Wiedau}
\institute{DEXPI e.V. // Frankfurt a.M., Germany}
\date{\today}

%----------------------------------------------------------------------------------------
%    PRESENTATION SLIDES
%----------------------------------------------------------------------------------------

\begin{document}
\begin{frame}[plain]
    \titlepage
\end{frame}

\begin{frame}{Overview}
    % Throughout your presentation, if you choose to use \section{} and \subsection{} commands, these will automatically be printed on this slide as an overview of your presentation
    \tableofcontents
\end{frame}

%------------------------------------------------
\section{Introduction}
%------------------------------------------------

\begin{frame}{Dr. Elena Marquez \textit{(example for vita-slide)}}
    \begin{columns}[c] % Centered vertically
        \begin{column}{0.35\textwidth}
            \centering
            \includegraphics[width=0.9\linewidth]{img/Business_women_portrait.jpg}
        \end{column}        
        \begin{column}{0.6\textwidth}
            \begin{description}
                \item[Education:] PhD in Process Systems Engineering, ETH Zurich; MSc in Chemical Engineering
                \item[Experience:] 15+ years in digital transformation \& data governance in the process industry (oil \& gas, chemicals, EPC)
                \item[Expertise:] Data standards (e.g., DEXPI, ISO 15926), knowledge graphs, asset lifecycle information management
                \item[Role:] Leads enterprise-wide data strategy, bridging engineering, IT, and AI-driven analytics for operational excellence
            \end{description}
        \end{column}
    \end{columns}
\end{frame}

%------------------------------------------------
\section{First Section}
%------------------------------------------------
\begin{frame}{The Idea of DEXPI}
\begin{block}{Idea}
    We work to create an open, neutral, and reliable data exchange standard for the process industry to establish future-proof digitalized collaboration.   
\end{block}
\end{frame}

%------------------------------------------------

\begin{frame}{Bullet Points}
    \begin{itemize}
        \item Lorem ipsum dolor sit amet, consectetur adipiscing elit
        \item Aliquam blandit faucibus nisi, sit amet dapibus enim tempus eu
        \item Nulla commodo, erat quis gravida posuere, elit lacus lobortis est, quis porttitor odio mauris at libero
        \item Nam cursus est eget velit posuere pellentesque
        \item Vestibulum faucibus velit a augue condimentum quis convallis nulla gravida
    \end{itemize}
\end{frame}

%------------------------------------------------

\begin{frame}{Blocks of Highlighted Text}
    In this slide, some important text will be \alert{highlighted} because it's important. Please, don't abuse it.

    \begin{block}{Block}
        Sample text
    \end{block}

    \begin{alertblock}{Alertblock}
        Sample text in red box
    \end{alertblock}

    \begin{examples}
        Sample text in green box. The title of the block is ``Examples".
    \end{examples}
\end{frame}

%------------------------------------------------

% Double Block Test
% Double Block Test
\begin{frame}{Double Block test}
\begin{columns}[T,onlytextwidth]
    \column{0.49\textwidth}
    \begin{block}{Format}
        \begin{itemize}
          \item Lorem ipsum dolor sit amet, consectetur adipiscing elit
          \begin{itemize}
            \item Sed do eiusmod tempor incididunt ut labore et dolore magna aliqua
          \end{itemize}
          \item Ut enim ad minim veniam, quis nostrud exercitation ullamco
          \item Duis aute irure dolor in reprehenderit in voluptate velit esse
        \end{itemize}
    \end{block}
    
    \column{0.49\textwidth}
    \begin{block}{Content}
        \begin{itemize}
          \item Excepteur sint occaecat cupidatat non proident
          \item Sunt in culpa qui officia deserunt mollit anim id est laborum
          \item Curabitur pretium tincidunt lacus, nulla gravida orci
        \end{itemize}
    \end{block}
\end{columns}
\end{frame}

%------------------------------------------------


\begin{frame}{Multiple Columns}
    \begin{columns}[c] % The "c" option specifies centered vertical alignment while the "t" option is used for top vertical alignment

        \column{.45\textwidth} % Left column and width
        \textbf{Heading}
        \begin{enumerate}
            \item Statement
            \item Explanation
            \item Example
        \end{enumerate}

        \column{.45\textwidth} % Right column and width
        Lorem ipsum dolor sit amet, consectetur adipiscing elit. Integer lectus nisl, ultricies in feugiat rutrum, porttitor sit amet augue. Aliquam ut tortor mauris. Sed volutpat ante purus, quis accumsan dolor.

    \end{columns}
\end{frame}

%------------------------------------------------
\section{Second Section}
%------------------------------------------------

\begin{frame}{Table}
    \begin{table}
        \begin{tabular}{l l l}
            \toprule
            \textbf{Treatments} & \textbf{Response 1} & \textbf{Response 2} \\
            \midrule
            Treatment 1         & 0.0003262           & 0.562               \\
            Treatment 2         & 0.0015681           & 0.910               \\
            Treatment 3         & 0.0009271           & 0.296               \\
            \bottomrule
        \end{tabular}
        \caption{Table caption}
    \end{table}
\end{frame}

%------------------------------------------------

\begin{frame}{Theorem}
    \begin{theorem}[Mass--energy equivalence]
        $E = mc^2$
    \end{theorem}
\end{frame}

%------------------------------------------------

\begin{frame}{Figure}
    Uncomment the code on this slide to include your own image from the same directory as the template .TeX file.
    %\begin{figure}
    %\includegraphics[width=0.8\linewidth]{test}
    %\end{figure}
\end{frame}

%------------------------------------------------

\begin{frame}[fragile] % Need to use the fragile option when verbatim is used in the slide
    \frametitle{Citation}
    An example of the \verb|\cite| command to cite within the presentation:\\~

    There are several DEXPI publications that can be cited: \cite{wiedau2019enpro}, \cite{oeing2022using}, \cite{cameron2024dexpi}, \cite{otten2024dexpi}.
\end{frame}

%------------------------------------------------

\begin{frame}{References}
    \footnotesize
    \bibliography{reference.bib}
    \bibliographystyle{apalike}
\end{frame}

%------------------------------------------------

\begin{frame}
    \Huge{\centerline{\textbf{The End}}}
\end{frame}


\end{document}